\section{Background and related work}\label{sec:related}
% Sources: docs/VERIFICHE-chiuse.md (verified: Fagin, Che, Fricker, UCP, Cliffhanger,
% FairRide, Zhang et al., Cloudflare taxonomy of 1 Jul 2026, IETF drafts) and
% docs/UNDERTOW-CONTESTO.md §7. \todo{verify} marks every source not yet verified on the
% primary: the proposals attributed to Zhang et al., the "unique URLs" framing of the
% Cloudflare blog, SemDN (bib entry "check"). Cloudflare blog of 2 Apr 2026 (Wildani and
% Ahmad): date, SIEVE/S3-FIFO "initial experiments" sentence and "separate cache layer for
% AI traffic" (long term) verified on the page, docs/PREREG-simulatore-fase2.md, amendment 1.
% Not in references.bib (no \cite until added and verified): Breakwater, DAGOR, Protego,
% TopFull, Denning 1968, Mattson et al. 1970, SHARDS, Counter Stacks.

\subsection{Background}\label{sec:background}
% Cloudflare taxonomy Search / Agent / Training (1 Jul 2026, cloudflare2026options,
% verified); the 15 Sep 2026 change is an announcement, for new domains, on ad pages only
% (section E: never "default block"). IETF webbotauth / aipref drafts (C6).
\todo{background}

\subsection{Related work}\label{sec:rw}
\paragraph{Automated traffic and web caches.}
Zhang et al.~\cite{zhang2025} vary the share of AI traffic from 0\% to 100\% on Varnish,
Apache Traffic Server and Memcached, and measure the overall hit ratio and that of human
traffic.
Their workloads fix the overlap between human and AI content, at about 20\% in the
Wikimedia setup and 10\% in the synthetic trace, and do not vary it.
A Cloudflare post~\cite{cloudflare2026cache} reports initial experiments in which SIEVE or
S3-FIFO in place of LRU could let human traffic keep the same hit rate with or without AI
traffic, and expects a separate cache layer for AI traffic to be the best way forward in the
long term.
Zhang et al. also discuss remedies such as newer eviction policies, learned caching and
admission control\todo{verify}.
We agree on the exhaustive class: varied on its own volume, it costs about 0.98 origin
requests per added request in our testbed.\claim{A7}
We add four things.
First, we measure the marginal cost at the origin and compare it across three coexisting
classes, rather than the hit ratio.\claim{A2}
Second, for the agentic class the sign of that cost depends on the width of the class's
working set (\cref{sec:r-sign}), and the agentic clients of our honeypot sit near the
narrow end by that measure.\claim{A1}\claim{C3}
Third, the overlap that these workloads fix changes the measured cost of a class
(\cref{sec:r-overlap}).\claim{A3}
Fourth, in our testbed the hit ratio of the exhaustive class rises from 0.179 to 0.200 as
the agentic share grows under the separated mapping, so isolating the agentic class in a
separate cache would remove a measured benefit (supported).\claim{D8}

\paragraph{Allocation, sharing and feedback.}
Utility-based cache partitioning~\cite{qureshi2006ucp} and Cliffhanger~\cite{cidon2016cliffhanger}
steer cache space by the gradient of the hit or miss ratio with respect to the memory
allocated to a workload.
We vary the composition of the traffic that shares a fixed cache, not the allocation.
FairRide~\cite{pu2016fairride} treats free-riding in a shared cache as a problem of
fairness and incentives; we treat the overlap between working sets as a bias in the
measured cost of a class.\claim{A3}
Admission control driven by feedback on load or latency is well established, for example
in Breakwater, DAGOR, Protego and TopFull \todo{add to references.bib and verify}.
Feedback control is not a contribution of this paper, and we do not propose a controller.

\paragraph{Working sets, characteristic time and miss-ratio curves.}
The working set of a program \todo{cite and verify: Denning 1968} is the concept we use
for a class of traffic.
The characteristic-time approximation of LRU is due to Fagin~\cite{fagin1977}, was
rediscovered and named by Che et al.~\cite{che2002}, and was formalised by Fricker et
al.~\cite{fricker2012}.
We use it to interpret the ordering of our results, and its quantitative prediction on
our testbed failed (\cref{sec:m-tc}).\claim{B3}\claim{B4}
Stack algorithms and their sampled approximations, such as SHARDS
\todo{cite and verify: Mattson et al. 1970, SHARDS, Counter Stacks}, estimate the miss
ratio as a function of cache size; we measure how the cost of one class changes with the
composition of the traffic and the width of its working set.

\paragraph{Delivery architectures for agents.}
SemDN~\cite{hua2026semdn}\todo{verify} proposes a delivery architecture that retrieves
and caches content at semantic granularity for LLM agents.
We do not propose a delivery architecture, and we have not tested our findings on
SemDN.
We ask a different question: whether, inside a shared infrastructure, the cost of a class
of traffic can be treated as a stable property of that class.\claim{D1}
